%%%%%%%%%%%%%%%%%%%%%%%%%%%%%%%%%%%%%%%%%%%%%%%%%%%%%%%%%%%%%%%%%%%%%%%%
% Part 1: Sections 1-4 (Introduction, Background, Contributions, Math)
%%%%%%%%%%%%%%%%%%%%%%%%%%%%%%%%%%%%%%%%%%%%%%%%%%%%%%%%%%%%%%%%%%%%%%%%

\section{Introduction}
\label{sec:introduction}

\subsection{Problem Statement}

The hedging of derivative instruments represents a fundamental challenge in quantitative finance. Traditional Black-Scholes delta hedging~\citep{black1973pricing} assumes continuous trading, zero transaction costs, and complete markets---assumptions that fail in practice. Deep hedging~\citep{buehler2019deep} reformulates this as a stochastic optimal control problem solved via neural network function approximation, naturally incorporating market frictions and risk preferences.

Despite significant progress, several fundamental challenges remain. First, \textbf{model uncertainty}: standard deep hedging trains on fixed Heston parameters; distributional shifts during crises degrade performance. Second, \textbf{training instability}: the two-stage CVaR$\to$entropic protocol~\citep{kozyra2019deep} exhibits abrupt objective transitions. Third, \textbf{regime dependence}: different architectures excel in different market conditions, yet single-model deployment is standard.

We address these challenges with three novel approaches:
\begin{enumerate}[noitemsep,leftmargin=*]
    \item \textbf{W-DRO-T:} Wasserstein DRO regularization for distributional robustness
    \item \textbf{3SCH:} Mixed-loss homotopy for smooth training transitions
    \item \textbf{RSE:} Regime-dependent gating over diverse base hedgers
\end{enumerate}

\subsection{Market Frictions in Real-World Financial Markets}

\begin{itemize}[noitemsep]
    \item \textbf{Transaction costs:} Real hedging incurs trading costs (SPY $\approx 3$ bps vs NIFTY $\approx 18$ bps including STT, stamp duty, and slippage).
    
    \item \textbf{Discrete rebalancing:} Hedging occurs at discrete intervals (e.g., daily rebalancing with $\sim 30$ steps for a 30-day option).
    
    \item \textbf{Model uncertainty:} Stochastic volatility models (e.g., Heston) face parameter instability and calibration errors across market regimes.
    
    \item \textbf{Incomplete markets:} Stochastic volatility prevents perfect replication of derivatives.
    
    \item \textbf{Position limits:} Regulatory and risk constraints restrict allowable hedge ratios and exposures.
    
    \item \textbf{India-specific frictions:} NIFTY/BANKNIFTY options involve STT (0.05\%), wider spreads, discrete strikes (50-point intervals), and SEBI position limits.
    
    \item \textbf{Strategic implication:} High trading frictions favor low-turnover hedging strategies over frequent rebalancing.
\end{itemize}

\subsection{Neural Network Representation of Hedging Policies}

The hedging strategy is parameterised as a neural network policy:
\begin{equation}
    \delta_k = \delta_{\max} \cdot \tanh(F_\theta(I_{0:k}, \delta_{k-1})), \qquad \delta_{\max} = 1.5,
\end{equation}
where $F_\theta$ is a neural network (LSTM, Transformer, or ensemble) and $I_k = (S_k/S_0, \log(S_k/K), \sqrt{v_k}, \tau_k, \Delta_k^{\mathrm{BS}}) \in \R^5$ is the feature vector. The $\tanh$ bounding ensures $|\delta_k| \leq 1.5$, preventing gradient explosions and enforcing economically sensible position limits. The gradient of the delta output:
\begin{equation}
    \frac{\partial \delta_k}{\partial h_k} = \delta_{\max}(1 - \tanh^2(W_\delta h_k + b_\delta)) \cdot W_\delta,
\end{equation}
is bounded by $\delta_{\max}\|W_\delta\|$ and vanishes only in the saturated regime.

\subsection{Motivation for Deep Learning in Hedging}

Deep hedging~\citep{buehler2019deep} has emerged as the dominant paradigm for learning optimal hedging policies. However, a fundamental vulnerability persists: models trained on paths from a \emph{fixed} stochastic volatility model degrade under distributional shift. LSTM hedgers trained on normal-regime Heston parameters exhibit CVaR$_{95}$ degradation exceeding $4\times$ under COVID-type volatility.

Our audit of baseline models reveals distinct regime-dependent strengths:
\begin{itemize}[noitemsep]
    \item \textbf{LSTM:} Best in stable, trending markets (lowest trading volume)
    \item \textbf{Transformer:} Best tail risk during volatile periods (attention captures clustering)
    \item \textbf{Signature models:} Best for path-dependent dynamics (captures path geometry)
\end{itemize}

%===============================================================
\section{Background and Literature Review}
\label{sec:background}

\subsection{Classical Hedging Theory}

Black--Scholes delta hedging prescribes $\delta_t^{\mathrm{BS}} = \Phi(d_1)$ where $d_1 = \frac{\log(S/K)+(r+\sigma^2/2)\tau}{\sigma\sqrt\tau}$. This assumes geometric Brownian motion, continuous trading, zero transaction costs, and market completeness. Under stochastic volatility (Heston model), the market is incomplete and hedging residuals persist.

\subsection{Deep Hedging Framework}

\citet{buehler2019deep} introduced neural-network-parameterised hedging under convex risk measures, demonstrating that learned policies outperform Black--Scholes delta hedging in the presence of transaction costs. \citet{kozyra2019deep} proposed two-stage curriculum training (CVaR pretraining followed by entropic fine-tuning) and introduced delta bounding via $\tanh$ for training stability. \citet{horvath2021deep} extended the framework to rough volatility models. \citet{carbonneau2021deep} applied deep hedging to long-term insurance derivatives.



\subsection{Attention Mechanisms and Transformer Architectures}

\citet{vaswani2017attention} introduced the Transformer; applications to financial time series include Informer~\citep{zhou2021informer} and SigFormer~\citep{tong2023sigformer}. For hedging, the causal self-attention mechanism captures long-range dependencies in market features that recurrent networks may miss. The causal mask enforces non-anticipativity: step $k$ cannot attend to future steps $k'>k$.

\textbf{Related DRO work.} \citet{blanchet2019quantifying} established Wasserstein-based DRO theory with tractable gradient-norm duality. \citet{esfahani2018data} proved finite-sample performance guarantees for data-driven DRO. \citet{lutkebohmert2019robust} showed robust approaches reduce hedging error during COVID-type volatility.

\textbf{Related RL work.} \citet{kolm2019dynamic} and \citet{cao2021deep} explored RL approaches. \citet{neagu2025drl} compared eight DRL algorithms. \citet{haarnoja2018soft} introduced SAC; \citet{huang2025sac} added CVaR constraints.

\textbf{Ensembles and regime switching.} \citet{dietterich2000ensemble} established ensemble variance reduction theory. \citet{hamilton1989new} introduced regime-switching models. Our RSE adapts the mixture-of-experts paradigm to hedging.

\subsection{Universal Approximation in Neural Networks}

LSTMs and Transformers are universal approximators for sequential mappings. Path signatures~\citep{lyons1998differential} provide universal features for path-dependent functionals---the signature uniquely determines the path up to tree-like equivalence.

%===============================================================
\section{Research Contributions}
\label{sec:contributions}

\subsection{Overview of Studied Architectures}

Seven architectures evaluated: LSTM, Transformer, W-DRO-T, 3SCH, RSE. All share identical data pipeline, feature engineering, and two-stage baseline protocol. The experimental framework: 10 independent seeds (42--942), bootstrap CIs (10,000 resamples), Holm--Bonferroni corrected paired tests, real-market validation on SPY and NIFTY across four crisis scenarios.

\subsection{Novel Approaches Introduced}

\begin{table}[H]
\centering\small
\caption{Novel approaches: design rationale.}
\begin{tabular}{@{}lllll@{}}
\toprule
\textbf{Model} & \textbf{Gap} & \textbf{Innovation} & \textbf{Theory} & \textbf{Ref.} \\\midrule
W-DRO-T & Distributional shift & Grad-norm penalty & Wasserstein DRO duality & \citep{blanchet2019quantifying} \\
3SCH & Training instability & Mixed-loss homotopy & Multiobjective Optimisation & \citep{kozyra2019deep} \\
RSE & Regime dependence & Gated ensemble & Ensemble variance reduction & Gap analysis \\

\bottomrule
\end{tabular}
\end{table}

\subsubsection{Wasserstein DRO Transformer (W-DRO-T)}

W-DRO-T optimises worst-case risk over a Wasserstein ball: $\theta^* = \argmin_\theta \sup_{\QQ: W_1(\QQ,\PP) \leq \varepsilon} \E_\QQ[\rho(\mathrm{P\&L})]$. Via duality~\citep{blanchet2019quantifying}, this reduces to gradient-norm regularisation. Three-phase training: CVaR pretraining $\to$ DRO-annealed entropic ($\varepsilon: 0 \to 0.1$) $\to$ stress testing. Achieves 46\% CVaR improvement during COVID on SPY and 41\% Basel~III/IV capital savings.

\subsubsection{Three-Stage Curriculum Hedger (3SCH)}

3SCH inserts a mixed-loss homotopy $\mathcal{L}(\alpha) = \alpha \cdot \cvar_{0.95} + (1-\alpha) \cdot \rho_\lambda$ with $\alpha: 0.8 \to 0.2$ between CVaR and entropic stages. Justified by Berge's maximum theorem ensuring upper hemicontinuous solution paths. Best NIFTY normal CVaR$_{95}$ of 622.26---optimal for high-cost emerging markets with trading volume of only 0.90.

\subsubsection{Regime-Switching Ensemble (RSE)}

RSE combines frozen LSTM, Transformer, and Signature hedgers via regime-dependent gating. Regime features (RV, trend, momentum, jump) drive a neural classifier mapping to $K=4$ regimes; a learnable affinity matrix $A \in \R^{4 \times 3}$ produces model weights. Best CVaR$_{95} = 3.109 \pm 0.010$ ($-3.29\%$ vs LSTM, $p < 0.0001$, Cohen's $d = -7.41$), lowest CV = 0.30\%.

%===============================================================
\section{Mathematical Framework}
\label{sec:math}

\subsection{Market Model}

\subsubsection{Heston Stochastic Volatility Model}

We work on $(\Omega, \F, (\F_t)_{t \in [0,T]}, \PP)$ satisfying the usual conditions (right-continuous, complete), supporting a two-dimensional Brownian motion $(W^S, W^v)$ with $\langle W^S, W^v \rangle_t = \rho\,dt$, $\rho \in (-1,0)$. The Heston model~\citep{heston1993closed}:
\begin{align}
    dS_t &= r S_t \, dt + \sqrt{v_t}\, S_t \, dW_t^S, \qquad S_0 > 0, \\
    dv_t &= \kappa(\theta_v - v_t) \, dt + \xi \sqrt{v_t} \, dW_t^v, \qquad v_0 > 0.
\end{align}

\begin{remark}[Feller Condition]
The variance process remains strictly positive when $2\kappa\theta_v > \xi^2$. Our calibration: $2 \times 1.0 \times 0.04 = 0.08 > 0.04 = \xi^2$.
\end{remark}

\subsubsection{Model Parameters}

$S_0 = K = 100$, $v_0 = 0.04$, $\kappa = 1.0$, $\theta_v = 0.04$, $\xi = 0.2$, $\rho = -0.7$, $r = 0.05$, $T = 30/365$, $N = 30$.

\subsection{Discretisation}

\subsubsection{Euler--Maruyama Scheme}

Full-truncation scheme on equidistant grid $\Delta t = T/N$:
\begin{align}
    S_{k+1} &= S_k \exp\!\bigl[(r - \tfrac{1}{2}v_k^+)\Delta t + \sqrt{v_k^+\Delta t}\;\epsilon_k^S\bigr], \\
    v_{k+1} &= v_k + \kappa(\theta_v - v_k)\Delta t + \xi\sqrt{v_k^+\Delta t}\;\epsilon_k^v,
\end{align}
where $v_k^+ = \max(v_k, 0)$ and $(\epsilon_k^S, \epsilon_k^v) \sim \mathcal{N}(0, \Sigma)$ with $\Sigma_{12} = \rho$.

\subsection{Deep Hedging Formulation}

\subsubsection{Hedging P\&L Definition}

For European call payoff $Z = (S_T - K)^+$ with $N = 30$ hedging steps:
\begin{equation}
    \mathrm{P\&L}(\delta) = -Z + \sum_{k=0}^{N-1} \delta_k (S_{t_{k+1}} - S_{t_k}) - c_{\mathrm{tc}} \sum_{k=0}^{N} |\delta_k - \delta_{k-1}| S_{t_k}
\end{equation}
with $c_{\mathrm{tc}} = 0.001$ and $\delta_{-1} = \delta_N = 0$ (flat entry/exit).

\subsubsection{Neural Policy Parameterisation}

$\theta^* = \argmin_\theta \rho(-\mathrm{P\&L}(\delta^\theta))$ with $\delta_k = \delta_{\max} \cdot \tanh(F_\theta(I_{0:k}, \delta_{k-1}))$, $\delta_{\max} = 1.5$.

\subsubsection{Feature Engineering}

$I_k = (S_k/S_0, \log(S_k/K), \sqrt{v_k}, \tau_k, \Delta_k^{\mathrm{BS}}) \in \R^5$ where $\tau_k = (T-t_k)/T$ is normalised time-to-maturity and $\Delta_k^{\mathrm{BS}}$ is Black--Scholes delta.

\subsection{Risk Measures}

\subsubsection{Conditional Value at Risk (CVaR)}

\begin{definition}[CVaR]
For loss $L = -\mathrm{P\&L}$ at level $\alpha \in (0,1)$:
\begin{align}
    \VaR_\alpha(L) &= \inf\{l \in \R : \PP(L \le l) \ge \alpha\}, \\
    \cvar_\alpha(L) &= \frac{1}{1-\alpha}\int_\alpha^1 \VaR_u(L)\,du = \E[L \mid L \geq \VaR_\alpha(L)].
\end{align}
CVaR is a coherent risk measure~\citep{artzner1999coherent}: monotonicity, translation invariance, positive homogeneity, and subadditivity.
\end{definition}

\begin{proposition}[Rockafellar--Uryasev Dual~\citep{rockafellar2000optimization}]
\begin{equation}
    \cvar_\alpha(L) = \inf_{\nu \in \R}\Bigl\{\nu + \frac{1}{1-\alpha}\,\E\bigl[(L-\nu)^+\bigr]\Bigr\}.
\end{equation}
Differentiable in $\theta$ (a.e.), enabling gradient-based optimisation.
\end{proposition}

\subsubsection{Entropic Risk Measure}

\begin{definition}[Entropic Risk]
With risk aversion $\lambda > 0$:
\begin{equation}
    \rho_\lambda(X) = \frac{1}{\lambda}\log\E[\exp(-\lambda X)].
\end{equation}
Dual representation~\citep{follmer2002convex}: $\rho_\lambda(X) = \sup_{\QQ \ll \PP}\{\E_\QQ[-X] - \frac{1}{\lambda}\mathrm{KL}(\QQ\|\PP)\}$.
\end{definition}

\begin{remark}[CVaR vs.\ Entropic: Qualitative Differences]
CVaR focuses exclusively on the tail ($\alpha$-quantile and beyond), producing strategies that aggressively protect against worst-case losses but may over-trade. Entropic risk weighs all outcomes exponentially, producing smoother strategies. The abrupt switch between these objectives motivates 3SCH.
\end{remark}

\subsubsection{Dual Representations of Risk Measures}

\begin{definition}[Generalised DRO Risk]
Given divergence $D(\QQ\|\PP)$ and radius $\varepsilon > 0$:
$\rho^D_\varepsilon(X) = \sup_{\QQ: D(\QQ\|\PP) \le \varepsilon} \E_\QQ[-X]$.
For $D = \mathrm{KL}$: entropic risk ($\varepsilon = 1/\lambda$). For $D = W_1$: Wasserstein DRO.
\end{definition}

\begin{remark}[Comparison of Ambiguity Sets]
\textbf{KL ball}: requires $\QQ \ll \PP$, adversary can only \emph{reweight} existing scenarios. \textbf{Wasserstein ball}: does \emph{not} require absolute continuity, adversary can \emph{move} mass to new regions. For crisis hedging, Wasserstein is more appropriate because distributional shifts involve \emph{new} volatility regimes (COVID $\sigma = 65\%$) not seen during training ($\sigma \approx 20\%$).
\end{remark}

\begin{definition}[$1$-Wasserstein Distance]
For probability measures $\mu, \nu$ on $(\R^d, \|\cdot\|_2)$:
\begin{equation}
    W_1(\mu,\nu) = \inf_{\gamma \in \Gamma(\mu,\nu)} \int \|x-y\|_2\,d\gamma(x,y).
\end{equation}
Kantorovich--Rubinstein duality: $W_1(\mu,\nu) = \sup_{\|f\|_{\mathrm{Lip}} \le 1}\{\int f\,d\mu - \int f\,d\nu\}$.
\end{definition}

